% Part: lambda-calculus
% Chapter: church-rosser
% Section: parallel-beta-reduction

\documentclass[../../../include/open-logic-section]{subfiles}

\begin{document}

\olfileid[id]{lam}{cr}{pb}

\olsection{Reduksi-$\beta$ Paralel}

Kita memperkenalkan gagasan \emph{reduksi-$\beta$ paralel} dan membuktikan
bahwa reduksi tersebut mempunyai sifat Church--Rosser.

\begin{defn}[Reduksi-$\beta$ paralel, $\bredpar$] \ollabel{defn:bredpar}
  Reduksi paralel ($\bredpar$) pada suku didefinisikan secara induktif sebagai
  berikut:
  \begin{enumerate}
    \item \ollabel{defn:bredpar1} $x \bredpar x$.
    \item \ollabel{defn:bredpar2} Jika $N \bredpar N'$, maka
      $\lambd[x][N] \bredpar \lambd[x][N']$.
    \item \ollabel{defn:bredpar3} Jika $P \bredpar P'$ dan $Q \bredpar Q'$,
      maka $PQ \bredpar P'Q'$.
    \item \ollabel{defn:bredpar4} Jika $N \bredpar N'$ dan $Q \bredpar Q'$,
      maka $(\lambd[x][N])Q \bredpar \Subst{N'}{Q'}{x}$.
  \end{enumerate}
\end{defn}

Reduksi-$\beta$ paralel memungkinkan kita mereduksi sebarang jumlah redex
dalam satu suku melalui satu langkah. Reduksi ini berbeda dari
reduksi-$\beta$ karena kita hanya dapat mengontraksi redex yang sudah terdapat
dalam suku asal, bukan redex yang baru muncul akibat reduksi-$\beta$ paralel.
Sebagai contoh, suku $(\lambd[f][fx])(\lambd[y][y])$ hanya dapat direduksi-$\beta$
secara paralel menjadi dirinya sendiri atau menjadi $(\lambd[y][y])x$, tetapi
tidak langsung lebih lanjut menjadi~$x$, meskipun suku itu tereduksi-$\beta$
menjadi~$x$; redex terakhir baru muncul sesudah satu langkah reduksi-$\beta$
paralel. Namun, langkah reduksi-$\beta$ paralel kedua menghasilkan~$x$.

\begin{thm}\ollabel{thm:refl}
  $M \bredpar M$.
\end{thm}

\begin{proof}
  Latihan.
\end{proof}

\begin{prob}
  Buktikan \olref[lam][cr][pb]{thm:refl}.
\end{prob}

\begin{defn}[Pengembangan lengkap-$\beta$]\ollabel{defn:bcd}
  \emph{Pengembangan lengkap-$\beta$}~$\bcd{M}$ dari $M$ didefinisikan secara
  induktif sebagai berikut:
  \begin{align}
    \bcd{x} &= x \ollabel{defn:bcd1} \\
    \bcd{(\lambd[x][N])} &= \lambd[x][\bcd{N}] \ollabel{defn:bcd2}\\
    \bcd{(PQ)} &= \bcd{P}\bcd{Q} && \text{jika $P$ bukan abstrak-$\lambd$}
    \ollabel{defn:bcd3} \\
    \bcd{((\lambd[x][N])Q)} &= \Subst{\bcd{N}}{\bcd{Q}}{x} \ollabel{defn:bcd4}
  \end{align}
\end{defn}

Sesuai namanya, pengembangan lengkap-$\beta$ suatu suku merupakan ``reduksi
paralel lengkap.'' Dalam reduksi-$\beta$ paralel kita masih dapat memilih untuk
tidak mengontraksi suatu redex; dalam pengembangan lengkap, semua redex harus
dikontraksi. Karena itu, pengembangan lengkap dari
$(\lambd[f][fx])(\lambd[y][y])$ adalah $(\lambd[y][y])x$, bukan suku itu
sendiri.

\begin{editorial}
  Definisi ini bermasalah karena cara mendefinisikan fungsi secara rekursif
  pada suku-$\lambd$ belum diperkenalkan. Hal ini akan diperbaiki kelak.
\end{editorial}

\begin{lem}\ollabel{lem:comp}
  Jika $M \bredpar M'$ dan $R \bredpar R'$, maka $\Subst{M}{R}{y}
  \bredpar \Subst{M'}{R'}{y}$.
\end{lem}

\begin{proof}
  Dengan induksi pada !!{derivation} $M \bredpar M'$.
  \begin{enumerate}
    \item Langkah terakhir adalah \olref{defn:bredpar1}: Latihan.
    \item Langkah terakhir adalah \olref{defn:bredpar2}: Maka
      $M=\lambd[x][N]$ dan $M'=\lambd[x][N']$, dengan $N\bredpar N'$.
      Kita ingin membuktikan bahwa
      $\Subst{(\lambd[x][N])}{R}{y}\bredpar
      \Subst{(\lambd[x][N'])}{R'}{y}$, yaitu
      $\lambd[x][\Subst{N}{R}{y}]\bredpar
      \lambd[x][\Subst{N'}{R'}{y}]$. Hal ini langsung mengikuti
      \olref{defn:bredpar2} dan hipotesis induksi.
    \item Langkah terakhir adalah \olref{defn:bredpar3}: Latihan.
    \item Langkah terakhir adalah \olref{defn:bredpar4}: $M$ adalah
      $(\lambd[x][N])Q$ dan $M'$ adalah $\Subst{N'}{Q'}{x}$. Kita ingin
      membuktikan bahwa
      $\Subst{((\lambd[x][N])Q)}{R}{y}\bredpar
      \Subst{\Subst{N'}{Q'}{x}}{R'}{y}$, yaitu
      $(\lambd[x][\Subst{N}{R}{y}])\Subst{Q}{R}{y}\bredpar
      \Subst{\Subst{N'}{R'}{y}}{\Subst{Q'}{R'}{y}}{x}$. Hal ini mengikuti
      \olref{defn:bredpar4} dan hipotesis induksi.
  \end{enumerate}
\end{proof}

\begin{prob}
  Lengkapi pembuktian \olref[lam][cr][pb]{lem:comp}.
\end{prob}

\begin{lem}\ollabel{lem:cont}
  Jika $M \bredpar M'$, maka $M' \bredpar \bcd{M}$.
\end{lem}
\begin{proof}
  Dengan induksi pada !!{derivation} $M \bredpar M'$.
  \begin{enumerate}
    \item Aturan terakhir adalah \olref{defn:bredpar1}: Latihan.
    \item Aturan terakhir adalah \olref{defn:bredpar2}: $M$ adalah
      $\lambd[x][N]$ dan $M'$ adalah $\lambd[x][N']$, dengan
      $N\bredpar N'$. Kita ingin menunjukkan
      $\lambd[x][N']\bredpar\bcd{(\lambd[x][N])}$, yaitu
      $\lambd[x][N']\bredpar\lambd[x][\bcd{N}]$ berdasarkan
      \olref{defn:bcd2}. Hal ini mengikuti \olref{defn:bredpar2} dan
      hipotesis induksi.
    \item Aturan terakhir adalah \olref{defn:bredpar3}: $M$ adalah $PQ$ dan
      $M'$ adalah $P'Q'$ untuk suatu $P$, $Q$, $P'$, dan~$Q'$, dengan
      $P\bredpar P'$ dan $Q\bredpar Q'$. Berdasarkan hipotesis induksi,
      $P'\bredpar\bcd{P}$ dan $Q'\bredpar\bcd{Q}$.
      \begin{enumerate}
        \item Jika $P=\lambd[x][N]$ untuk suatu $x$ dan~$N$, maka
          $P'=\lambd[x][N']$ untuk suatu $N'$ dengan $N\bredpar N'$.
          Berdasarkan hipotesis induksi, $N'\bredpar\bcd{N}$ dan
          $Q'\bredpar\bcd{Q}$. Maka
          $(\lambd[x][N'])Q'\bredpar\Subst{\bcd{N}}{\bcd{Q}}{x}$ berdasarkan
          \olref{defn:bredpar4}.
        \item Jika $P$ bukan abstrak-$\lambd$, maka
          $P'Q'\bredpar\bcd{P}\bcd{Q}$ berdasarkan
          \olref{defn:bredpar3}; ruas kanannya adalah $\bcd{PQ}$ berdasarkan
          \olref{defn:bcd3}.
      \end{enumerate}
    \item Aturan terakhir adalah \olref{defn:bredpar4}: $M$ adalah
      $(\lambd[x][N])Q$ dan $M'$ adalah $\Subst{N'}{Q'}{x}$ untuk suatu
      $x$, $N$, $Q$, $N'$, dan~$Q'$, dengan $N\bredpar N'$ dan
      $Q\bredpar Q'$. Berdasarkan hipotesis induksi,
      $N'\bredpar\bcd{N}$ dan $Q'\bredpar\bcd{Q}$. Berdasarkan
      \olref{lem:comp},
      $\Subst{N'}{Q'}{x}\bredpar\Subst{\bcd{N}}{\bcd{Q}}{x}$; ruas kanan
      ini tepat sama dengan $\bcd{((\lambd[x][N])Q)}$.
  \end{enumerate}
\end{proof}

\begin{prob}
  Lengkapi pembuktian \olref[lam][cr][pb]{lem:cont}.
\end{prob}

\begin{thm}\ollabel{thm:cr}
  $\bredpar$ mempunyai sifat Church--Rosser.
\end{thm}

\begin{proof}
  Langsung dari \olref{lem:cont}.
\end{proof}

\end{document}
